\begin{table}[t]
    \centering
    \small
    \resizebox{\ifdim\width>\textwidth\textwidth\else\width\fi}{!}{%
\begin{tabular}{@{}llc@{}}
        \toprule
        \textbf{Input type} & \textbf{Conversion rule} & \textbf{PET effects} \\
        \midrule
        \texttt{pearson\_r}, \texttt{partial\_r}, \texttt{r\_from\_partial\_eta2}
            & used directly as $r$ (with $r = \sqrt{\eta_p^2}$ for a 1-df contrast) & 20 \\
        \texttt{kendall\_tau}
            & $r = \sin(\pi\tau/2)$ (Greiner's relation) & 1 \\
        \texttt{unstd\_beta\_*}
            & test-statistic recovery, \eqref{eq:recovery} & 12 \\
        \texttt{narrative\_null}
            & imputed $z = 0$, $v = 1/(n-3)$ & 7 \\
        \midrule
        \multicolumn{2}{@{}l}{\textbf{Total usable amyloid-PET effects}} & \textbf{40} \\
        \bottomrule
    \end{tabular}%
}
    \caption{Effect harmonisation rules, fixed in advance and applied uniformly. All effects are
    pooled as Fisher $z$ with $v = 1/(n-3)$. Test-statistic recovery raises the usable
    amyloid-PET set from 21 native correlations to 40 effects and is what brings the two largest
    cohorts onto the common metric.}
    \label{tab:conversion}
\end{table}
